
Comme @@M0@@ est @@M1@@-régulier, @@M2@@ par 1.4; d'où @@M3@@. Or, @@M4@@ par 1.8 (i) puisque @@M5@@ est un @@M6@@-faisceau par 2.5. Par suite
@@M7@@
Enfin, 2.7.2 donne (ii), et 2.7.1 et 2.7.2 donnent (i).

\textbf{Lemme 2.8.} \emph{On suppose que @@M8@@ et @@M9@@. Soient @@M10@@ et @@M11@@. Alors} :
@@M12@@

En effet, soit @@M13@@; alors il résulte en comparant des polynômes que @@M14@@, que @@M15@@ et que @@M16@@. Soit @@M17@@; alors @@M18@@ et par 2.3, @@M19@@, ce qui égale @@M20@@. Donc @@M21@@ et comme @@M22@@, on a @@M23@@ par 2.7 (ii); d'où l'assertion.

\textbf{Lemme 2.9.} \emph{On suppose que @@M24@@, que @@M25@@ et que @@M26@@ soit normal. Alors} :
\begin{enumerate}[label=(\roman*)]
\item \emph{@@M27@@ pour @@M28@@,}
\item \emph{@@M29@@ pour tout @@M30@@,}
\item \emph{@@M31@@ pour @@M32@@,}
\end{enumerate}
\emph{où @@M33@@ et @@M34@@ sont des polynômes universels à coefficients dans @@M35@@, qui ne dépendent ni de @@M36@@ ni de @@M37@@.}

